\documentclass[11pt,a4paper]{article}
\usepackage[margin=27mm]{geometry}
\usepackage{fontspec}
\setmainfont{lmroman10-regular.otf}[BoldFont=lmroman10-bold.otf,
  ItalicFont=lmroman10-italic.otf,BoldItalicFont=lmroman10-bolditalic.otf]
\usepackage{amsmath,amssymb,amsthm}
\usepackage{microtype}
\usepackage{xurl}
\usepackage[unicode,colorlinks=true,linkcolor=blue,urlcolor=blue]{hyperref}
\usepackage{bookmark}
\setcounter{secnumdepth}{0}
\setlength{\emergencystretch}{2em}
\providecommand{\tightlist}{\setlength{\itemsep}{0pt}\setlength{\parskip}{0pt}}
\title{Uniform integer representation counts and reflected spans\\
for quadratic Bose--Chowla sets}
\author{Sage Kim\\Peking University}
\date{Draft for external mathematical review\\5 November 2026}
\hypersetup{pdfauthor={Sage Kim},pdftitle={Uniform integer representation counts and reflected spans for quadratic Bose-Chowla sets},pdfsubject={External-review draft; no novelty claim or general Erdos 864 resolution}}
\begin{document}
\maketitle
\begin{abstract}

We consider the complete quadratic Bose–Chowla sets with arbitrary
primitive elements and arbitrary positive representatives in one period.
We prove a uniform absolute local limit for the ordered integer
representation function of $b_1+b_2+b_3-b_4$. The proof parametrizes
the modular solution fibre, bounds its entire omitted part by $6p$,
and uses a degree-five Dirichlet character together with an exact
integer-carry selector. We then deduce that the least inclusive span
of an admissible reflected union is $(3+o(1))p^2$, even when every
integer reflection height is allowed. Both conclusions persist after
arbitrary $o(p)$ deletions. The reflection and deletion results are
elementary consequences of the representation estimate and endpoint
occupancy; they are not separate character-sum advances.

\end{abstract}

\section{Introduction}\label{introduction}

How much can one shorten a reflected Sidon construction without creating a second repeated sum? This question is part of the study of sumsets of Sidon sequences developed by \href{https://www.sciencedirect.com/science/article/pii/0022314X9190083N}{Erdős and Freud}. In \href{https://hrj.episciences.org/125/pdf}{Erdős's 1992 account, §2, pp.39--40}, the one-exception sum problem is stated as follows: among sets of integers \(A\subseteq\{1,\ldots,N\}\) for which at most one sum has more than one unordered representation, including diagonals, determine the largest possible cardinality. That account records the attainable leading constant \(2/\sqrt3\) and proposes its asymptotic optimality.

The elementary reflected construction starts with a dense Sidon set \(B\) and forms \(B\cup(S-B)\). A sufficiently large height \(S\) separates its sum families, leaving only the central sum with repeated representations. For a complete quadratic Bose--Chowla base of size \(p\) and diameter asymptotic to \(p^2\), a feasible choice has inclusive hull length \((3+o(1))p^2\) and cardinality asymptotic to \(2p\). Could another reflection height produce a smaller leading span? Could deleting \(o(p)\) base elements first make such a shortening possible?

The obstruction is an integer representation problem. When the two reflected pieces are separated, an internal/mixed sum collision has the form \(b_1+b_2+b_3-b_4=S\). Solving this equation modulo \(p^2-1\) is insufficient: a modular witness can lie at a different integer lift \(S+k(p^2-1)\). One needs to count witnesses at the prescribed integer height, with its actual positive representatives.

The analytic result of this paper is Theorem 1, a uniform absolute asymptotic formula for those integer counts, for every primitive element, every one-period cut, and every integer target. Its proof combines an explicit rational parametrization, a uniform bound on the entire discarded modular fibre, and an exact integer-carry selector. The only external analytic input is a proved polynomial Dirichlet-character estimate. Corollary 1.1 then gives deletion stability and optimal reflected span: even with arbitrary integer heights, overlapping pieces and \(o(p)\) deletions, this construction family retains the leading constant \(2/\sqrt3\). These are applications of the representation estimate, not separate advances in character-sum theory.

The restriction to complete quadratic bases and their sublinear deletions is essential to the scope of the result. We do not prove an upper bound for arbitrary admissible sets in the general Erdős problem \#864, and we make no claim of literature priority.

\section{1. Definitions and main results}\label{definitions-and-main-results}

Throughout, \(p\) tends to infinity through odd primes,

\[K=\mathbb F_{p^2},\qquad M=p^2-1,\]

and \(\theta\) is a generator of \(K^{\times}\). For \(t\in\mathbb Z/M\mathbb Z\) define

\[\beta(a)=\operatorname{rep}_{[1,M]}
       \bigl(\log_\theta(\theta+a)+t\bigr),\qquad
B=B_{p,\theta,t}=\{\beta(a):a\in\mathbb F_p\}.
\tag{1}\]

Here residue zero has representative \(M\). Write \(\mathcal B_p\) for the family of all these bases. For a finite integer set \(C\), let

\[R_C(s)=\#\{(a,b,c,d)\in C^4:a+b+c-d=s\},\qquad s\in\mathbb Z.
\tag{2}\]

Coordinates are ordered and all repetitions are permitted. Separately, call an integer set admissible if at most one integer sum has more than one \textbf{unordered} representation \(a\le b\), including diagonal pairs. For nonempty \(C\) and every integer \(S\), put

\[A_C(S)=C\cup(S-C),\qquad
\ell_C(S)=\max A_C(S)-\min A_C(S)+1,\]

\[L^*(C)=\min\{\ell_C(S):S\in\mathbb Z, A_C(S)\text{ admissible}\}.
\tag{3}\]

The minimization uses the actual inclusive hull length, allowing overall integer translation. In particular it imposes neither separation nor positivity on \(S\). Existence of a feasible \(S\) is proved below.

Let \(U_1,U_2,U_3\) be independent uniform random variables on \((0,1)\), and define

\[F(y)=\mathbb P\{y<U_1+U_2+U_3<y+1\}.
\tag{4}\]

\textbf{Theorem 1 (uniform integer representation counts).}

With the definitions above,

\[\varepsilon_p:=\sup_{B\in\mathcal B_p}\sup_{s\in\mathbb Z}
\left|\frac{R_B(s)}{p^2}-F\!\left(\frac{s}{M}\right)\right|
\longrightarrow0.
\tag{5}\]

The supremum includes all integer targets, with absolute rather than relative error. In particular no positive lower bound is asserted arbitrarily close to a zero of the limiting profile at a prescribed rate.

\textbf{Corollary 1.1 (deletion stability and reflected-span applications).}

\textbf{(i)} For every \(B\in\mathcal B_p\), every \(D\subseteq B\), and every integer \(s\),

\[0\le R_B(s)-R_{B\setminus D}(s)\le8p|D|.
\tag{6}\]

Consequently, if \(k_p\) is any nonnegative integer sequence with \(k_p/p\to0\), then

\[\sup_{\substack{B\in\mathcal B_p, D\subseteq B\\|D|\le k_p}}
\sup_{s\in\mathbb Z}
\left|\frac{R_{B\setminus D}(s)}{p^2}-F\!\left(\frac{s}{M}\right)\right|
\le\varepsilon_p+\frac{8k_p}{p}\longrightarrow0.
\tag{7}\]

\textbf{(ii)} For every such sequence \(k_p\), the survivors \(C=B\setminus D\) are nonempty for large \(p\), \(L^*(C)\) exists, and

\[\sup_{\substack{B\in\mathcal B_p, D\subseteq B\\|D|\le k_p}}
\left|\frac{L^*(B\setminus D)}{p^2}-3\right|\longrightarrow0.
\tag{8}\]

Taking \(D=\varnothing\) gives the complete-base all-height result. All limits are absolute, uniform over primitive elements, cuts and stated deletions. No relative error near a zero of \(F\), explicit rate, or explicit prime threshold is asserted. Arbitrary Sidon bases and arbitrary lifts outside a single period are not included in (5), (7), or (8).

\section{2. Supporting lemmas}\label{supporting-lemmas}

\subsection{External analytic input}\label{external-analytic-input}

For a nonprincipal Dirichlet character \(\Xi\) modulo \(Q\in\mathbb F_p[T]\), extended by zero on nonunits, the proved function-field character bound gives

\[\left|\sum_{b\in\mathbb F_p}\Xi(T+b)\right|
\le(\deg Q-1)\sqrt p,\qquad \deg Q\ge2.
\tag{9}\]

This is the degree-one case of \href{https://arxiv.org/pdf/1204.0708v3}{Keating--Rudnick, §3.4, (3.31)--(3.33)}. The von Mangoldt weight of each monic linear polynomial is one. The bound requires nonprincipality, not primitivity. It rests on the proved function-field Riemann hypothesis and is the sole non-elementary analytic input here. It is not a conjectural appeal to number-field GRH.

\subsection{Lemma 2.1 (Sidon uniqueness and fixed-coordinate loads)}\label{lemma-2.1-sidon-uniqueness-and-fixed-coordinate-loads}

The map \(a\mapsto\beta(a)\) is injective, and \(B\) is Sidon modulo \(M\), and hence over the integers, with diagonal sums included. Any nonzero difference has at most one ordered representation. For a fixed modular target \(s\), at most \(2p\) modular witnesses to \(b_1+b_2+b_3-b_4\equiv s\pmod M\) contain a specified element in a specified slot.

\textbf{Proof.} Put \(L(a)=\theta+a\). These \(p\) field elements are distinct and nonzero. Equality of two modular pair sums is equivalent to

\[L(a)L(b)=L(c)L(d).\]

After cancelling \(\theta^2\) and comparing coefficients in the basis \(1,\theta\), we get \(a+b=c+d\) and \(ab=cd\). Thus the unordered pairs are the roots, with multiplicity, of the same quadratic. Difference uniqueness follows by rewriting an equality of differences as an equality of pair sums; the swapped alternative forces difference zero.

To bound a positive-slot load, fix that element and fix the negative element in \(p\) ways. The other two positive elements have a prescribed modular pair sum, with at most two ordered realizations. For a fixed negative slot, fix any positive slot first and use the same argument. This proves the \(2p\) bound for each of the four slots. \hfill\(\square\)

\subsection{Lemma 2.2 (parametrization and the entire exceptional complement)}\label{lemma-2.2-parametrization-and-the-entire-exceptional-complement}

Fix \(s\in\mathbb Z\), put \(\gamma=\theta^{s-2t}\), and let

\[\mathcal T_s=\{(a,b,c,e)\in\mathbb F_p^4:
                  L(a)L(b)L(c)=\gamma L(e)\}.\]

There is a subset \(\mathcal G_s\subseteq\mathcal T_s\) satisfying

\[|\mathcal G_s|=(p-e_s)(p-1),\quad 0\le e_s\le4,\qquad
|\mathcal T_s\setminus\mathcal G_s|\le6p.
\tag{10}\]

In particular the total number of singular witnesses described below is at most \(2p\), rather than merely \(O(p)\) per row.

\textbf{Proof.} Write conjugation as \(x\mapsto\bar x=x^p\) and let \(\Delta=\theta-\bar\theta\ne0\). The exponent \(s-2t\) is forced by the translation contribution \(t+t+t-t=2t\). Fix \(a\), set

\[k=L(a)/\gamma,\qquad z=kL(b).\]

When \(z\ne\bar z\), the equation \(zL(c)=L(e)\) and its conjugate give the unique solution

\[L(c)=\frac{\Delta(1-\bar z)}{z-\bar z}.
\tag{11}\]

For the displayed value \(C_0\), direct conjugation shows \(C_0-\bar C_0=\Delta\) and \(zC_0-\bar z\bar C_0=\Delta\). Thus \(C_0-\theta\) and \(zC_0-\theta\) belong to \(\mathbb F_p\), supplying unique \(c,e\). Both \(C_0\) and \(zC_0\) are nonzero because their conjugate differences are nonzero.

Put \(w=\bar\gamma/(\bar\theta+a)-\bar\theta\). Discard a row \(a\) if \(k\in\mathbb F_p\), \(w\in\mathbb F_p\), or \(w\in\{-\theta,-\bar\theta\}\). At most four rows are discarded:

\begin{itemize}
\item
  Two distinct rows with real \(k\) would, by subtraction, imply \(1/\gamma\in\mathbb F_p\), which makes real \(L(a)/\gamma\) impossible. Hence at most one row has real \(k\).
\item
  The equation \(w=\bar w\), after multiplication by nonzero denominators, is

  \[\bar\gamma(\theta+a)-\gamma(\bar\theta+a)
       +\Delta(\theta+a)(\bar\theta+a)=0.\]

  This polynomial in \(a\) has degree two and leading coefficient \(\Delta\ne0\), so there are at most two such rows.
\item
  The equation \(w=-\theta\) is \(\bar\gamma=-\Delta(\bar\theta+a)\), at most one row. The equation \(w=-\bar\theta\) would give \(\gamma=0\) and never occurs.
\end{itemize}

On a retained row \(k\ne\bar k\), and the unique root of \(z-\bar z\) as a function of \(b\) is

\[r=-\frac{k\theta-\bar k\bar\theta}{k-\bar k}\in\mathbb F_p.\]

The ratio is fixed by conjugation. For \(b\ne r\), (11) becomes

\[L(c)=\lambda\frac{b-w}{b-r},\qquad
\lambda=-\frac{\Delta\bar k}{k-\bar k}\ne0.
\tag{12}\]

Define \(\mathcal G_s\) to contain exactly these solutions on retained rows with \(b\ne r\). They have distinct first two coordinates, and each retained row has \(p-1\) of them. This proves its exact size.

It remains to count ALL omitted witnesses. If \(z=\bar z\), taking conjugate differences in \(zL(c)=L(e)\) forces \(z\Delta=\Delta\). Thus no solution exists for real \(z\ne1\). For \(z=1\) there are exactly \(p\) solutions \(c=e\). Such parameter pairs satisfy

\[L(a)L(b)=\gamma.\]

The pair-product uniqueness already proved in Lemma 2.1 gives at most two ordered \((a,b)\) for this fixed \(\gamma\), in the entire parameter plane. Therefore all singular witnesses together number at most \(2p\).

Every nonsingular parameter pair has exactly one continuation by (11). Those omitted from \(\mathcal G_s\) are in at most four discarded rows, and contribute at most \(4p\). All other nonsingular witnesses are retained. Thus the full complement has size at most \(4p+2p=6p\). This counts continuations, not just exceptional parameter pairs. \hfill\(\square\)

\subsection{Lemma 2.3 (joint Fourier control and endpoint occupancy)}\label{lemma-2.3-joint-fourier-control-and-endpoint-occupancy}

Give each \((a,b,c,e)\in\mathcal G_s\) mass \(p^{-2}\) at \((\beta(a)/M,\beta(b)/M,\beta(c)/M)\), obtaining a measure \(\mu_s\) on \(\mathbb T^3\). Its total mass is \(1+O(1/p)\). At a mode \((i,j,h)\) nonzero modulo \(M\), its unnormalized Fourier sum is at most

\[\begin{cases}
4p\sqrt p,&(j,h)\not\equiv(0,0)\pmod M,\\
(p-1)(\sqrt p+4),&j\equiv h\equiv0\pmod M.
\end{cases}
\tag{13}\]

All constants are uniform in \(\theta,t,s\). Hence these measures converge uniformly against every fixed continuous torus function and every fixed box with null boundary to Lebesgue measure. For the bases, each fixed interval \(I\subset(0,1)\) satisfies

\[|B\cap MI|=(|I|+o(1))p
\tag{14}\]

uniformly over \(B\in\mathcal B_p\). In particular

\[\min B/M\to0,\qquad\max B/M\to1.
\tag{15}\]

\textbf{Proof.} Choose \(\chi:K^{\times}\to\mathbb C^{\times}\) of exact order \(M\), with \(\chi(\theta)=\exp(2\pi i/M)\). Translation affects a Fourier coefficient only by a constant phase. On a retained row, (12) shows that the varying part of the coefficient is

\[\sum_{b\ne r}\chi^j(b+\theta)\chi^h(b-w)\chi^{-h}(b-r).
\tag{16}\]

For nonreal \(\alpha\), let \(m_\alpha(T)\) be its minimal polynomial over \(\mathbb F_p\), and put

\[Q(T)=m_\theta(T)m_{-w}(T)(T+r).
\tag{17}\]

The two quadratic factors are irreducible. They are distinct because \(w\ne-\theta,-\bar\theta\). The linear factor differs from both. Therefore \(Q\) is squarefree of degree five and the evaluation map is

\[\mathbb F_p[T]/(Q)\ \simeq\ K\times K\times\mathbb F_p,
\qquad f\longmapsto(f(\theta),f(-w),f(-r)).\]

On its units define the Dirichlet character

\[\Xi(f)=\chi^j(f(\theta))\chi^h(f(-w))\chi^{-h}(f(-r)),\]

restricting the last factor to \(\mathbb F_p^{\times}\), and extend it by zero on nonunits. Evaluation of \(T+b\) gives exactly (16), including a zero at \(b=r\), even if the last local character is principal. If \((j,h)\not\equiv(0,0)\), at least one full \(K^{\times}\) component is nonprincipal. The Chinese remainder theorem lets that component vary independently, so \(\Xi\) is nonprincipal. Cancellation after restriction to the base field does not make this character principal. Equation (9) gives \(4\sqrt p\) per row, proving the first bound in (13).

If \(j=h=0\) modulo \(M\), then \(i\ne0\) and the degree-two modulus \(m_\theta\), with character \(f\mapsto\chi^i(f(\theta))\), gives

\[\left|\sum_a\chi^i(L(a))\right|\le\sqrt p.\]

Removing at most four rows changes this sum by at most four; each retained row has exactly \(p-1\) columns. This proves the second bound.

For any fixed nonzero integer Fourier mode, its residue is nonzero for large \(p\), so division by \(p^2\) in (13) gives a uniformly vanishing coefficient. The zero coefficient tends uniformly to one by (10). Uniform trigonometric approximation of a fixed continuous torus function, using finitely many fixed modes, proves uniform weak convergence. A fixed box can be bounded above and below by continuous functions whose integrals differ arbitrarily little, because its boundary has Lebesgue measure zero. This gives the box assertion, including boxes meeting the representative seam at 1.

The degree-two bound also gives \(O(p^{-1/2})\) for each nonzero fixed mode of the base measure \(p^{-1}\sum_{b\in B}\delta_{b/M}\). The same approximation proves (14). In particular every fixed endpoint interval has positive asymptotic occupancy, proving (15). Alternatively, Lemma 2.1 bounds each fixed coordinate atom in \(\mu_s\) by \(2/p\); thus the representative seam cannot hold positive limiting mass. \hfill\(\square\)

\subsection{Lemma 2.4 (uniform moving slabs)}\label{lemma-2.4-uniform-moving-slabs}

Let \(\mu\) be any finite positive measure on \((0,1]^3\). Partition that cube into \(H^3\) half-open boxes of side \(1/H\), and put

\[\delta_H(\mu)=\max_Q|\mu(Q)-|Q||,
\qquad E_y=\{x\in(0,1]^3:0<x_1+x_2+x_3-y\le1\}.\]

For every real \(y\),

\[|\mu(E_y)-|E_y||\le H^3\delta_H(\mu)+12/H.
\tag{18}\]

\textbf{Proof.} The union of boxes wholly contained in \(E_y\) bounds it from below, and the union of boxes meeting \(E_y\) bounds it from above. Any box in the difference has closure meeting a plane with coordinate sum \(y\) or \(y+1\). Since the coordinate sum varies by at most \(3/H\) within a box, the difference lies in the two bands

\[|x_1+x_2+x_3-y|\le3/H,\qquad
|x_1+x_2+x_3-(y+1)|\le3/H.\]

Each band has volume at most \(6/H\): fix \(x_1,x_2\) and integrate the possible interval for \(x_3\). Thus their union has volume at most \(12/H\), independently of \(y\). Summing discrepancies over at most \(H^3\) boxes in either approximation gives (18). Half-open endpoints and coordinate seams do not affect these pointwise inclusions. \hfill\(\square\)

\subsection{Lemma 2.5 (elementary deletion and reflection bounds)}\label{lemma-2.5-elementary-deletion-and-reflection-bounds}

Let \(C\) be any finite integer Sidon set of size \(n\ge1\), and \(Z\subseteq C\). Then for every integer \(s\),

\[0\le R_C(s)-R_{C\setminus Z}(s)\le8n|Z|.
\tag{19}\]

For every integer \(S\),

\[|C\cap(S-C)|\le2,
\qquad A_C(S)\text{ admissible }\Longrightarrow R_C(S)\le6n.
\tag{20}\]

\textbf{Proof.} The fixed-coordinate proof in Lemma 2.1 applies over the integers too: fixing any one slot and a suitable second slot leaves a prescribed pair sum with at most two ordered realizations. Thus each deleted element lies in at most \(2n\) witnesses in each specified slot. A union bound over four slots and \(|Z|\) elements gives (19). It permits repeated coordinates and empty survivors.

The intersection \(C\cap(S-C)\) consists of the endpoints of a pair of sum \(S\) inside \(C\), so has at most two elements. The reflected union is invariant under \(x\mapsto S-x\). Any repeated sum \(u\ne S\) would reflect to a different repeated sum \(2S-u\). Therefore, in an admissible union, only the centre sum \(S\) may repeat.

For every witness \(a+b+c-d=S\), compare the unordered pairs \(\{a,b\}\) and \(\{d,S-c\}\) in the union. Unless their sum is \(S\) or the two pairs coincide, this is a forbidden off-centre collision. All permitted witnesses therefore lie in the union of the following three classes:

\[a+b=S,\ c=d;\qquad
b+c=S,\ a=d;\qquad
a+c=S,\ b=d.\]

Each class has at most two ordered choices for its displayed sum pair and \(n\) choices for its equal coordinates. Summing gives \(6n\). The classes need not be disjoint; overcounting only strengthens the upper bound. Diagonal pair sums are covered. \hfill\(\square\)

\subsection{Lemma 2.6 (exact hull and a feasible height)}\label{lemma-2.6-exact-hull-and-a-feasible-height}

For nonempty integer Sidon \(C\), let \(m=\min C\), \(v=\max C\), and \(d=v-m\). At every integer height,

\[\ell_C(S)=d+|S-(m+v)|+1.
\tag{21}\]

The height \(S_0=3v-m+1>2v\) is feasible and

\[L^*(C)\le3d+2.
\tag{22}\]

\textbf{Proof.} The hull endpoints are \(\max(v,S-m)\) and \(\min(m,S-v)\). Splitting at \(S=m+v\) yields (21). At \(S_0\) the two pieces are disjoint. Their internal sum families \(C+C\) and \(2S_0-(C+C)\) are unique by Sidonicity, lying below and above \(S_0\), respectively. Mixed sums have form \(S_0+x-y\) and are unique away from \(S_0\) by difference uniqueness. A collision between a low internal sum and a mixed sum would force

\[S_0=a+b+c-d'\le3v-m,\]

which is impossible. A high/mixed collision reflects to such a low/mixed one. The only possible repeated sum is \(S_0\), so the union is admissible and has length \(3d+2\). The feasible lengths are a nonempty set of positive integers, hence have a minimum. \hfill\(\square\)

\section{3. Proof of Theorem 1 and Corollary 1.1}\label{proof-of-theorem-1-and-corollary-1.1}

\subsection{3.1. Proof of Theorem 1: exact integer lift and uniform error}\label{proof-of-theorem-1-exact-integer-lift-and-uniform-error}

Use the modular fibre and its retained measure from Lemmas 2.2--2.3. Put \(y=s/M\). For a modular witness whose first three representatives are \(b_1,b_2,b_3\), the fourth representative is the unique element of \(\{1,\ldots,M\}\) congruent to \(b_1+b_2+b_3-s\). Therefore it satisfies the INTEGER equation if and only if

\[0<\frac{b_1+b_2+b_3-s}{M}\le1.
\tag{23}\]

The weak upper endpoint selects representative \(M\); the lower endpoint 0 must be excluded. Thus, with \(E_y\) as in Lemma 2.4,

\[0\le\frac{R_B(s)}{p^2}-\mu_s(E_{s/M})\le\frac6p.
\tag{24}\]

This uses the full complement bound (10), including every large singular continuation. It would not follow merely by bounding the number of excluded parameter pairs.

Fix a positive integer \(H\). By Lemma 2.3 and finiteness of the grid,

\[\delta_{p,H}:=\sup_{\theta,t,s}\delta_H(\mu_s)\longrightarrow0
\quad\text{as }p\to\infty\text{ with }H\text{ fixed}.\]

Since \(|E_y|=F(y)\), Lemma 2.4 and (24) give the explicit qualitative error certificate

\[\varepsilon_p\le\frac6p+H^3\delta_{p,H}+\frac{12}{H}.
\tag{25}\]

Every term on the right is independent of the integer target. First let \(p\to\infty\) at fixed \(H\); then let \(H\to\infty\). This proves (5) simultaneously for all integers, including targets with \(s/M\) approaching the support endpoints. For targets outside \([-M,3M]\), both the integer count and profile vanish in any event. There is no unproved exchange with growing frequencies or shrinking boxes. \hfill\(\square\)

\subsection{3.2. Corollary 1.1(i): uniform deletion loss}\label{corollary-1.1i-uniform-deletion-loss}

Apply Lemma 2.5 with \(C=B\), \(n=p\), \(Z=D\). This proves (6). Combining its one-sided loss with both sides of (5) proves (7). No randomness or independence of the deletions and the target is assumed. \hfill\(\square\)

\subsection{3.3. Corollary 1.1(ii): every reflection height and survivor endpoints}\label{corollary-1.1ii-every-reflection-height-and-survivor-endpoints}

Let \(|D|\le k_p=o(p)\) and \(C=B\setminus D\). For each fixed \(0<\eta<1\), Lemma 2.3 places \((\eta+o(1))p\) parent elements in each of \((0,\eta M)\) and \(((1-\eta)M,M)\). Fewer than \(o(p)\) are deleted, so each interval still has survivors for large \(p\), uniformly in all choices. Consequently, writing \(m_C=\min C\), \(v_C=\max C\), \(d_C=v_C-m_C\),

\[m_C/M\to0,\qquad v_C/M\to1,\qquad d_C/M\to1.
\tag{26}\]

The profile \(F\) is continuous: after fixing the first two coordinates, moving \(y\) by \(h\) changes the allowable interval for the third coordinate by symmetric-difference length at most \(2|h|\), so \(|F(y+h)-F(y)|\le2|h|\). It is strictly positive on \((-1,3)\). Indeed a sum interval \((y,y+1)\) intersects \((0,3)\) precisely for such \(y\), and a small open cube inside \((0,1)^3\) then has all coordinate sums in that interval. Thus every compact nonempty \(J\subset(-1,3)\) has \(\min_J F>0\). By (7), uniformly for all \(S/M\in J\), the count \(R_C(S)\) has a positive constant times \(p^2\) witnesses eventually. Lemma 2.5 bounds it by \(6|C|\le6p\) at every admissible height. Hence no such compact-interior height is admissible for large \(p\).

By Lemma 2.6 and (26), \(L^*(C)\le3d_C+2=(3+o(1))M\). If, for fixed \(0<\epsilon<1\), a feasible height had \(\ell_C(S)\le(3-\epsilon)M\) along arbitrarily large primes, the exact formula (21) would give

\[\left|\frac S M-\frac{m_C+v_C}{M}\right|
\le2-\epsilon+o(1).\]

Using (26), eventually

\[S/M\in[-1+\epsilon/2,\,3-\epsilon/2].\]

This fixed compact interval has just been excluded. Thus \(L^*(C)\ge(3-o(1))M\) uniformly. Since \(M/p^2\to1\), (8) follows. This argument includes overlapping and negative-height reflections. \hfill\(\square\)

\section{4. Corollaries and finite limitations}\label{corollaries-and-finite-limitations}

\textbf{Corollary 4.1 (explicit profile and interior coverage).} Put

\[G(z)=\frac{z_+^3-3(z-1)_+^3+3(z-2)_+^3-(z-3)_+^3}{6}.\]

Here \(x_+=\max(x,0)\). Then \(F(y)=G(y+1)-G(y)\). It is zero outside \([-1,3]\), positive inside, and satisfies \(F(y)=F(2-y)\). In particular \(F(0)=F(2)=1/6\), \(F(1)=2/3\). For every compact nonempty \(J\subset(-1,3)\), all sufficiently large primes satisfy

\[R_{B\setminus D}(s)\ge\tfrac12(\min_J F)p^2\]

uniformly for \(s/M\in J\) and \(|D|\le k_p=o(p)\).

\textbf{Proof.} Inclusion-exclusion for the simplex \(x_1+x_2+x_3\le z\) in the unit cube gives \(G(z)\); subtract the two cumulative volumes. The substitution \(x_i\mapsto1-x_i\) gives the symmetry. Coverage follows from (7) and the positive minimum proved above. \hfill\(\square\)

\textbf{Corollary 4.2 (separated optimization and normalized size).} The minimum restricted to \(S>2\max C\) also equals \((3+o(1))p^2\) uniformly for the survivors in Corollary 1.1. Moreover

\[\sup_{S:\ A_C(S)\,\mathrm{admissible}}
\frac{|A_C(S)|}{\sqrt{\ell_C(S)}}\longrightarrow\frac2{\sqrt3}\]

uniformly over those survivors, with all integer heights allowed.

\textbf{Proof.} The unrestricted minimum is a lower bound for the separated minimum, while \(S_0\) in Lemma 2.6 is separated and has asymptotic length \(3p^2\). For every height, (20) gives \(2|C|-2\le|A_C(S)|\le2|C|\). Here \(|C|/p\to1\), and Corollary 1.1 bounds all feasible lengths below by \((3-o(1))p^2\); \(S_0\) supplies the matching asymptotic ratio. \hfill\(\square\)

\textbf{Corollary 4.3 (necessary deletion cost).} For every \(D\subseteq B\) with nonempty \(C=B\setminus D\), if \(A_C(S)\) is admissible then

\[|D|\ge\frac{R_B(S)-6|C|}{8p}
\tag{27}\]

whenever the numerator is positive. For any sequence of these choices with \(S/M\to y\in(-1,3)\),

\[\liminf\frac{|D|}{p}\ge F(y)/8.\]

\textbf{Proof.} Combine (6) and (20), then use (5) and continuity of \(F\). The term \(6|C|/p^2\) tends to zero because \(|C|\le p\). \hfill\(\square\)

This is a necessary cost, not a sharp cover theorem or stability under positive-proportion deletion. An optional admissible central element changes the reflected cardinality by at most one and cannot remove any existing collision, so it does not alter Corollary 4.2.

The following small examples prevent stronger finite misinterpretations. In \(\mathbb F_3[X]/(X^2-2)\), take \(\theta=1+X\), \(t=0\), so \(B=\{1,6,7\}\). Then \(R_B(-7)=0\) although \(F(-7/8)=1/3072>0\). At \(S=2\) the union \(\{-5,-4,1,6,7\}\) is admissible with \(R_B(2)=7\): its sole repeated sum is 2. Thus \(R_B(S)=0\) is not a necessary admissibility criterion for overlapping reflections. For the same base the unrestricted finite optimum is 7 at \(S=8\), whereas the separated finite optimum is 15 at \(S=16\). Neither (5) nor (8) asserts an unqualified finite equality.

\section{5. Relation to prior results}\label{relation-to-prior-results}

The Bose--Chowla construction and its Sidon proof are classical; \href{https://arxiv.org/pdf/2104.12711v3}{Nathanson, Theorem 1} gives a modern reproduction. Our analytic input (9) is already a consequence of proved function-field character theory. The application here requires the constrained parametrization, independent CRT factors, total exceptional-complement count, and the positive integer selector. Once compact-interior coverage and endpoint occupancy are available, the reflection obstruction, hull optimization and deletion robustness follow by the elementary arguments above. Full equality in (5) is stronger than the positivity needed for those consequences.

\href{https://www.numdam.org/item/10.5802/jtnb.1239.pdf}{Ortega--Prendiville, Theorems 1.2 and 6.3} provide ordinary Fourier discrepancy; their Theorem 1.15 concerns homogeneous equations with at least five variables. Neither is a direct statement of the prescribed four-variable integer local limit. Their proof's use of fourth moments in §5 explains the variable count. There is also a quantitative issue with an immediate four-factor transfer. For \(f=1_B\), \(g=(p/M)1_{[1,M]}\), and \(\Delta=\|\widehat f-\widehat g\|_\infty\), telescoping and Cauchy--Schwarz give

\[\sup_s|R_f(s)-R_g(s)|=O(\Delta M^{3/4}).
\tag{28}\]

Here for a real function \(h\) of finite support, \(R_h(s)=\sum_{a+b+c-d=s}h(a)h(b)h(c)h(d)\), and the Fourier transform is \(\widehat h(\alpha)=\sum_n h(n)\exp(2\pi i n\alpha)\) on \(\mathbb R/\mathbb Z\). Indeed \(\|\widehat f\|_2^2=p\), \(\|\widehat f\|_4^4=2p^2-p\), while \(\|\widehat g\|_2=O(1)\) and \(\|\widehat g\|_4^4=O(M)\); each remaining three-factor integral has size \(O(M^{3/4})\). The improved published bound \(\Delta=O(M^{3/8})\) supplies only \(O(M^{9/8})\) in (28), instead of the required \(o(p^2)=o(M)\). Furthermore

\[\Delta\ge\|\widehat f-\widehat g\|_2
       =\sqrt{p-p^2/M}\sim M^{1/4}.\]

Thus this particular absolute-value certificate cannot produce an \(o(M)\) error even at the best discrepancy order. This does not exclude other Fourier arguments. Lemma 2.3 uses joint Fourier control of constrained witnesses, not just a one-coordinate discrepancy.

\href{https://opikhurko.warwick.ac.uk/E/Pikhurko06dm.pdf}{Pikhurko}, §2 (7), §6 Lemma 10 and §7, gives interval/residue distribution, macroscopic triple-sum densities and randomized reflected-intersection averages. These checked statements do not supply the uniform prescribed integer count. \href{https://arxiv.org/pdf/2610.36207v1}{Carter--O'Bryant--Tao, Theorems 2.3--2.5} give diameter, position and gap estimates, which concern endpoint information rather than joint constrained fibres.

\href{https://arxiv.org/pdf/2301.12878v3}{Forey--Fresán--Kowalski}, introduction, §3 and §5, gives symmetric Sidon terminology and the algebraic setting of Bose sets, with a discussion of exponential-sum equidistribution. The latter varies character parameters and is not itself a statement about positive integer representatives at every target. \href{https://arxiv.org/pdf/2606.15041v2}{Ding, Theorems 1 and 4} addresses weighted one-coordinate distribution on arithmetic progressions. No deduction of (5) from these checked statements alone was established.

These comparisons identify the input and the extra reduction; they do not establish literature priority. In particular an elementary consequence or special application of a known character estimate must not be advertised as a new general character-sum theorem. Older original sources and broader arithmetic Fourier-transform treatments have not been exhaustively reviewed. There is no novelty claim here.

\section{6. Computational supplement}\label{computational-supplement}

The \href{https://github.com/miyosep/erdos-864-c021/blob/main/manuscript/notes/c021_editorial_20261105/COMPUTATIONAL_SUPPLEMENT.md}{computational supplement} provides the precise finite scopes, code, receipts and review limitations. The proofs use no experimental premise. Computations do not establish the asymptotic theorem, an explicit prime threshold or novelty. Theorem 1 and Corollary 1.1 concern the specified quadratic construction family and its sublinear deletions; the general Erdős \#864 upper-bound problem is outside their conclusion.

\sloppy\section{References}\label{references}

\begin{enumerate}
\def\labelenumi{\arabic{enumi}.}
\tightlist
\item
  J. P. Keating and Z. Rudnick, \emph{The variance of the number of prime polynomials in short intervals and in residue classes}, \href{https://arxiv.org/pdf/1204.0708v3}{arXiv:1204.0708v3}, §3.4, (3.31)--(3.33), pp.9--10. Exact external estimate used.
\item
  M. B. Nathanson, \emph{The Bose-Chowla argument for Sidon sets}, \href{https://arxiv.org/pdf/2104.12711v3}{arXiv:2104.12711v3}, J. Number Theory 238 (2022), 133--146, §2 Theorem 1.
\item
  M. Ortega and S. Prendiville, \emph{Extremal Sidon Sets are Fourier Uniform, with Applications to Partition Regularity}, \href{https://www.numdam.org/item/10.5802/jtnb.1239.pdf}{JTNB 35 (2023), 115--134}, Theorems 1.2, 1.15, 6.3 and §5 (5.4).
\item
  O. Pikhurko, \emph{Dense edge-magic graphs and thin additive bases}, \href{https://opikhurko.warwick.ac.uk/E/Pikhurko06dm.pdf}{Discrete Mathematics 306 (2006), 2097--2107}, §2 proof of Theorem 4, §6 Lemma 10 and §7.
\item
  D. Carter, K. O'Bryant and T. Tao, \emph{On the diameter and distribution of Sidon sets}, \href{https://arxiv.org/pdf/2610.36207v1}{arXiv:2610.36207v1}, Theorems 2.3--2.5.
\item
  A. Forey, J. Fresán and E. Kowalski, \emph{Sidon sets in algebraic geometry}, \href{https://arxiv.org/pdf/2301.12878v3}{arXiv:2301.12878v3}, introduction, Remark 2, §§3 and 5; published in IMRN 2024, no.8, 6400--6421.
\item
  Y. Ding, \emph{Dense finite Sidon sets on arithmetic progressions}, \href{https://arxiv.org/pdf/2606.15041v2}{arXiv:2606.15041v2}, Theorems 1 and 4, §2.1. The v2 title is used here.
\item
  P. Erdős and R. Freud, \emph{On sums of a Sidon-sequence}, \href{https://www.sciencedirect.com/science/article/pii/0022314X9190083N}{Journal of Number Theory 38 (1991), no.2, 196--205}, DOI 10.1016/0022-314X(91)90083-N.
\item
  P. Erdős, \emph{Some of my forgotten problems in number theory}, \href{https://hrj.episciences.org/125/pdf}{Hardy-Ramanujan Journal 15 (1992), 34--50}, DOI 10.46298/hrj.1992.125, especially §2, pp.39--40.
\end{enumerate}

\end{document}
